\section{2013 Algebra \& Number Theory}

\subsection*{Individual}

\begin{exercise}
\begin{enumerate}
    \item (15 pt) Classify finite groups of order 26 up to isomorphisms.
    \item (5 pt) For each finite group $G$ of order 26, describe the group $\mathrm{Aut}(G)$ of automorphisms of $G$.
\end{enumerate}
\end{exercise}

\begin{solution}
\textbf{Prerequisites / Core Concepts:}
\begin{itemize}
    \item \textbf{Sylow Theorems:} Used to determine the number and normality of subgroups of prime power order.
    \item \textbf{Semi-direct Products:} A way to construct a new group from two given groups, generalizing the direct product, particularly useful when one subgroup is normal.
    \item \textbf{Automorphism Groups of Cyclic Groups:} $\text{Aut}(\mathbb{Z}/n\mathbb{Z}) \cong (\mathbb{Z}/n\mathbb{Z})^\times$.
    \item \textbf{Chinese Remainder Theorem:} Used to decompose cyclic groups and their automorphism groups.
\end{itemize}

\textbf{Step 1: Classifying the groups of order 26 using Sylow Theorems.}

Let's start by looking at the order of our group, which is $|G| = 26$. We can prime factorize this as $2 \times 13$.

To understand the structure of $G$, we should examine its Sylow subgroups. Let $n_{13}$ be the number of Sylow 13-subgroups in $G$. By the Third Sylow Theorem, we know two things about $n_{13}$:
1. $n_{13} \equiv 1 \pmod{13}$
2. $n_{13}$ divides $2$ (the other prime factor of the group's order)

The only positive integer that satisfies both of these conditions is $n_{13} = 1$. What does this mean intuitively? Since there is only one Sylow 13-subgroup, let's call it $P_{13}$, it must be invariant under conjugation by any element in $G$. Therefore, $P_{13}$ is a \textit{normal} subgroup of $G$. Since its order is the prime $13$, $P_{13}$ must be cyclic, so $P_{13} \cong \mathbb{Z}/13\mathbb{Z}$.

Similarly, there is a Sylow 2-subgroup $P_2$ of order $2$, so $P_2 \cong \mathbb{Z}/2\mathbb{Z}$. The intersection $P_{13} \cap P_2$ must be just the identity element because their orders are coprime.

\textbf{Step 2: Constructing the groups via semi-direct products.}

Since $P_{13}$ is normal, $P_2 \cap P_{13} = \{e\}$, and $|P_{13}||P_2| = 13 \times 2 = 26 = |G|$, the group $G$ can be written as a semi-direct product of these two subgroups: $G \cong P_{13} \rtimes_\phi P_2 \cong \mathbb{Z}/13\mathbb{Z} \rtimes_\phi \mathbb{Z}/2\mathbb{Z}$.

The structure of this semi-direct product is entirely determined by a homomorphism $\phi: \mathbb{Z}/2\mathbb{Z} \to \text{Aut}(\mathbb{Z}/13\mathbb{Z})$.

Let's figure out what $\text{Aut}(\mathbb{Z}/13\mathbb{Z})$ looks like. For any prime $p$, the automorphism group of $\mathbb{Z}/p\mathbb{Z}$ is isomorphic to the multiplicative group $(\mathbb{Z}/p\mathbb{Z})^\times$, which is itself a cyclic group of order $p-1$. So, $\text{Aut}(\mathbb{Z}/13\mathbb{Z}) \cong (\mathbb{Z}/13\mathbb{Z})^\times \cong \mathbb{Z}/12\mathbb{Z}$.

Now we need to find all possible homomorphisms $\phi: \mathbb{Z}/2\mathbb{Z} \to \mathbb{Z}/12\mathbb{Z}$. Such a homomorphism is completely determined by where it sends the generator $1 \in \mathbb{Z}/2\mathbb{Z}$. Because the order of $1$ is $2$, the order of its image $\phi(1)$ must divide $2$. In the cyclic group $\mathbb{Z}/12\mathbb{Z}$, there are only two elements whose orders divide $2$: the identity element $0$ and the element $6$.

This gives us two possible maps:
1. \textbf{The trivial map:} $\phi(1) = 0$ (which corresponds to the identity automorphism $x \mapsto x$ in $\text{Aut}(\mathbb{Z}/13\mathbb{Z})$). The semi-direct product becomes a simple direct product, yielding the abelian group $\mathbb{Z}/13\mathbb{Z} \times \mathbb{Z}/2\mathbb{Z} \cong \mathbb{Z}/26\mathbb{Z}$ (the cyclic group $C_{26}$).
2. \textbf{The inversion map:} $\phi(1) = 6$ (which corresponds to the unique element of order 2 in $(\mathbb{Z}/13\mathbb{Z})^\times$, namely the inversion automorphism $x \mapsto -x$). This means the non-identity element of $\mathbb{Z}/2\mathbb{Z}$ acts on $\mathbb{Z}/13\mathbb{Z}$ by sending every element to its inverse. This defines the dihedral group $D_{26}$, which represents the symmetries of a regular 13-gon.

Thus, up to isomorphism, the only groups of order 26 are $C_{26}$ and $D_{26}$.

\textbf{Step 3: Finding the automorphism group of $C_{26}$.}

Now we need to describe $\text{Aut}(G)$ for each case. First, let $G = \mathbb{Z}/26\mathbb{Z}$.

Using the property of automorphisms of cyclic groups, we know $\text{Aut}(\mathbb{Z}/n\mathbb{Z}) \cong (\mathbb{Z}/n\mathbb{Z})^\times$.
So, $\text{Aut}(\mathbb{Z}/26\mathbb{Z}) \cong (\mathbb{Z}/26\mathbb{Z})^\times$.

By the Chinese Remainder Theorem, since $26 = 2 \times 13$ and $\gcd(2, 13) = 1$, we have:
$(\mathbb{Z}/26\mathbb{Z})^\times \cong (\mathbb{Z}/2\mathbb{Z})^\times \times (\mathbb{Z}/13\mathbb{Z})^\times$.

Since $(\mathbb{Z}/2\mathbb{Z})^\times \cong \{1\}$ (the trivial group) and $(\mathbb{Z}/13\mathbb{Z})^\times \cong \mathbb{Z}/12\mathbb{Z}$, we get:
$\text{Aut}(C_{26}) \cong \{1\} \times \mathbb{Z}/12\mathbb{Z} \cong \mathbb{Z}/12\mathbb{Z}$.

So the automorphism group is simply a cyclic group of order 12.

\textbf{Step 4: Finding the automorphism group of $D_{26}$.}

Next, let $G = $Dihedral group $D_{26}$. We can write its presentation as $D_{26} = \langle r, s \mid r^{13}=1, s^2=1, srs=r^{-1} \rangle$.

An automorphism $\sigma$ of $D_{26}$ must map generators to elements that satisfy the same relations.
The element $r$ has order 13. Its image $\sigma(r)$ must also have order 13. The only elements of order 13 in $D_{26}$ are the non-trivial powers of $r$. So, $\sigma(r) = r^k$ for some $k \in \{1, 2, \dots, 12\}$. There are $12$ choices for $k$.

The element $s$ has order 2, and it's not in the normal subgroup $\langle r \rangle$. Its image $\sigma(s)$ must also have order 2 and cannot be a power of $r$. The elements of order 2 in $D_{26}$ that are outside $\langle r \rangle$ are exactly the reflections, which have the form $r^j s$ for $j \in \{0, 1, \dots, 12\}$. There are $13$ choices for $j$.

Any choice of $k$ and $j$ uniquely defines an automorphism. So there are $12 \times 13 = 156$ automorphisms in total.
The group $\text{Aut}(D_{26})$ is isomorphic to the affine group $\text{Aff}(\mathbb{F}_{13}) \cong \mathbb{F}_{13} \rtimes \mathbb{F}_{13}^\times$.
\end{solution}

\begin{exercise}
Consider $f \in \mathbb{Z}_{>0}$ and nonzero vector spaces $V_i$ indexed by $i \in \mathbb{Z}/f\mathbb{Z}$. Suppose that there are linear maps $\phi_i: V_i \to V_{i+1}$ and $\psi_i: V_i \to V_{i-1}$ such that
\[ \phi_{i-1} \circ \psi_i = 0, \quad \psi_{i+1} \circ \phi_i = 0. \]
Prove that there exists lines $\ell_i \subset V_i$ for every $i \in \mathbb{Z}/f\mathbb{Z}$ such that
\[ \phi_i(\ell_i) \subset \ell_{i+1}, \quad \psi_i(\ell_i) \subset \ell_{i-1} \]
under one of the following two conditions:
\begin{enumerate}
    \item all $\psi_i = 0$, or
    \item $\dim V_i$ are equal to each other.
\end{enumerate}
\end{exercise}

\begin{solution}
\textbf{Prerequisites / Core Concepts:}
\begin{itemize}
    \item \textbf{Eigenvectors:} For a linear map $T: V \to V$, a line $\ell$ is invariant if $T(\ell) \subset \ell$.
    \item \textbf{Kernel and Image:} The "Orpheus condition" $\phi \psi = 0$ means the image of one map is inside the kernel of the other.
\end{itemize}

\textbf{(1) Step 1: Solving part (1) where all $\psi_i = 0$.}
Since all backward maps $\psi_i$ are zero, the condition $\psi_i(\ell_i) \subset \ell_{i-1}$ is always satisfied (since $0 \subset \ell_{i-1}$).
We only need to find lines $\ell_i$ such that $\phi_i(\ell_i) \subset \ell_{i+1}$.
Consider the composition map going around the circle:
\[ \Phi = \phi_f \circ \phi_{f-1} \circ \dots \circ \phi_1 : V_1 \to V_1 \]
Since we are over an algebraically closed field like $\mathbb{C}$, $\Phi$ has at least one eigenvector $v_1 \in V_1$.
Define $\ell_1 = \text{span}\{v_1\}$. Then define $\ell_2 = \text{span}\{\phi_1(v_1)\}$, and so on. If $\phi_i(v_i)$ is zero, we pick an arbitrary line in $V_{i+1}$.
By construction, $\phi_i(\ell_i) \subset \ell_{i+1}$ for all $i$.

\textbf{(2) Step 2: Solving part (2) with equal dimensions.}
When dimensions are equal, say $n$, and we have the relations $\phi \psi = 0$ and $\psi \phi = 0$.
If any $\phi_i$ or $\psi_i$ is zero, we can reduce the problem to part (1) or a simpler chain.
If all are non-zero, the Orpheus condition implies the space $V_i$ decomposes based on the kernels.
The problem can be solved using a fixed-point argument on the projective spaces. The space $X = \mathbb{P}(V_1) \times \dots \times \mathbb{P}(V_f)$ is compact. The relations define a closed subset of $X$ that is preserved by the "shift" operators.
The equal dimension ensures that the maps between projective spaces are well-behaved, and the algebraic set of such sequences of lines is non-empty.
\end{solution}

\begin{exercise}
For a parameter $t = (t_0, t_1, \dots, t_5) \in \mathbb{F}_5^6$ with $t_0 \neq 0$ and $\{t_i, i > 0\}$ an ordering of elements in $\mathbb{F}_5$, define a polynomial
\[ P_t(x) = (x - t_1)(x - t_2)(x - t_3) + t_0(x - t_4)(x - t_5). \]
\begin{enumerate}
    \item Show that $P_t(x)$ is irreducible in $\mathbb{F}_5[x]$;
    \item Show that two parameters $t, t'$ give the same polynomial if and only if $t_0 = t_0'$ and $\{t_4, t_5\} = \{t_4', t_5'\}$.
    \item Show that every irreducible cubic monic polynomial over $\mathbb{F}_5$ is obtained in this way.
\end{enumerate}
\end{exercise}

\begin{solution}
\textbf{Prerequisites / Core Concepts:}
\begin{itemize}
    \item \textbf{Irreducibility:} A cubic polynomial is irreducible over a field if and only if it has no roots in that field.
    \item \textbf{Counting Argument:} There are exactly 40 monic irreducible cubic polynomials over $\mathbb{F}_5$.
\end{itemize}

\textbf{Part (1): Prove irreducibility by checking roots.}
Since the degree is 3, we just check if $P_t(c) = 0$ for any $c \in \mathbb{F}_5$.
The set $\{t_1, \dots, t_5\}$ is a permutation of all 5 elements of $\mathbb{F}_5$.
- If $c \in \{t_1, t_2, t_3\}$, then $(c-t_1)(c-t_2)(c-t_3) = 0$. 
  Then $P_t(c) = 0 + t_0(c-t_4)(c-t_5)$. Since $c \neq t_4$ and $c \neq t_5$ (distinct elements) and $t_0 \neq 0$, the result is non-zero.
- If $c \in \{t_4, t_5\}$, then $t_0(c-t_4)(c-t_5) = 0$.
  Then $P_t(c) = (c-t_1)(c-t_2)(c-t_3) + 0$. Since $c$ is not one of $\{t_1, t_2, t_3\}$, the result is non-zero.
Since no element is a root, the polynomial is irreducible.

\textbf{Part (2): Identify unique parameters.}
The polynomial is monic of degree 3. It is determined by its values.
The sets $\{t_1, t_2, t_3\}$ and $\{t_4, t_5\}$ partition $\mathbb{F}_5$.
Two parameters give the same polynomial if and only if they result in the same coefficients.
The count of distinct sets $\{t_4, t_5\}$ is $\binom{5}{2} = 10$.
The number of choices for $t_0$ is 4.
Total constructed polynomials = $10 \times 4 = 40$.

\textbf{Part (3): Completeness.}
The number of monic irreducible cubics over $\mathbb{F}_q$ is $\frac{1}{3}(q^3 - q)$. For $q=5$:
\[ N_5(3) = \frac{125 - 5}{3} = 40. \]
Since we constructed 40 distinct polynomials and all are irreducible (part 1), we have found them all.
\end{solution}

\begin{exercise}
For $k$ a non-negative integer, let $V_k := \mathbb{R}[x]_{\leq k}$ be the vector space of real polynomials of degree at most $k$ with an action by $SL_2(\mathbb{R})$.
\end{exercise}

\begin{solution}
\textbf{Prerequisites / Core Concepts:}
\begin{itemize}
    \item \textbf{Representation of $SL_2$:} The group of $2 \times 2$ matrices acts on polynomials via fractional linear transformations.
    \item \textbf{Weight Spaces:} The diagonal matrices in $SL_2$ scale coordinates, identifying the "weights" of the representation.
    \item \textbf{Clebsch-Gordan Formula:} $V_m \otimes V_n \cong V_{m+n} \oplus V_{m+n-2} \oplus \dots \oplus V_{|m-n|}$.
\end{itemize}

\textbf{Part (1): Irreducibility.}
The representation $V_k$ has dimension $k+1$. The basis vectors $1, x, x^2, \dots, x^k$ are eigenvectors for the diagonal matrix $H = \begin{pmatrix} \lambda & 0 \\ 0 & \lambda^{-1} \end{pmatrix}$.
The eigenvalues (weights) are $\lambda^k, \lambda^{k-2}, \dots, \lambda^{-k}$. 
Since all weights are distinct and form a single string with step 2, the representation is irreducible.

\textbf{Part (2): Decomposition of Tensor Product.}
The map $\phi: V_m \otimes V_n \to V_{m+n}$ defined by $P(x, y) \mapsto P(x, x)$ is a surjective homomorphism.
Its kernel consists of polynomials that vanish when $y=x$, so they are divisible by $(y-x)$.
As shown in the original solution, multiplication by $(y-x)$ is an equivariant map from $V_{m-1, n-1}$ to $V_{m, n}$.
Iterating this process peels off one irreducible component at a time, resulting in the direct sum decomposition.

\textbf{Part (3): Invariants.}
The space of invariants $(V_\ell \otimes V_m \otimes V_n)^{SL_2}$ is the number of times the trivial representation $V_0$ appears in the tensor product.
By the Clebsch-Gordan formula, $V_\ell$ must appear in the decomposition of $V_m \otimes V_n$.
This occurs if and only if $|m-n| \leq \ell \leq m+n$ and the parity matches ($\ell \equiv m+n \pmod 2$).
These conditions are equivalent to the triangle inequality and even sum.
\end{solution}

\subsection*{Team}
\begin{exercise}
Let $A$ be an $n \times n$ skew symmetric real matrix.
\begin{enumerate}
    \item Prove that all eigenvalues are imaginary or zero and $e^A$ is orthogonal.
    \item Find conditions for $B = e^A$.
\end{enumerate}
\end{exercise}

\begin{solution}
\textbf{Prerequisites / Core Concepts:}
\begin{itemize}
    \item \textbf{Skew-Symmetric Matrix:} $A^T = -A$.
    \item \textbf{Orthogonal Matrix:} $B^T B = I$, which means it preserves the norm.
    \item \textbf{Matrix Exponential:} $e^A = I + A + A^2/2! + \dots$.
\end{itemize}

\textbf{Part (1): Eigenvalues and Orthogonality.}
- \textbf{Eigenvalues:} Let $\lambda$ be an eigenvalue with eigenvector $v$. Then $Av = \lambda v$.
  Consider the inner product $\langle v, Av \rangle = \lambda \|v\|^2$.
  Also $\langle v, Av \rangle = \langle A^T v, v \rangle = \langle -Av, v \rangle = -\bar{\lambda} \|v\|^2$.
  Thus $\lambda = -\bar{\lambda}$, meaning the real part of $\lambda$ is zero. So $\lambda$ is purely imaginary or zero.
- \textbf{Orthogonality:} $(e^A)^T = e^{A^T} = e^{-A}$.
  Then $(e^A)^T e^A = e^{-A} e^A = e^{A-A} = e^0 = I$. Thus $e^A$ is orthogonal.

\textbf{Part (2): Solvability.}
For an orthogonal matrix $B$ to be the exponential of a skew-symmetric matrix, it must have determinant 1 ($\det e^A = e^{\text{tr } A} = e^0 = 1$). 
The map $\exp: \mathfrak{so}(n) \to SO(n)$ is surjective.
Thus, $B$ must be in the \textbf{Special Orthogonal Group} $SO(n)$.
\end{solution}

\begin{exercise}
Let $SL_2(\mathbb{F}_p)$ be the group of $2 \times 2$ matrices over $\mathbb{F}_p$ with determinant 1.
\end{exercise}

\begin{solution}
\textbf{Prerequisites / Core Concepts:}
\begin{itemize}
    \item \textbf{Order of $GL_2$:} $(p^2-1)(p^2-p)$.
    \item \textbf{Characteristic Polynomial:} For a $2 \times 2$ matrix, it is $x^2 - \text{tr}(A)x + \det(A) = 0$.
\end{itemize}

\textbf{Part (1): Order of the group.}
The order of $GL_2(\mathbb{F}_p)$ is $(p^2-1)(p^2-p)$.
$SL_2(\mathbb{F}_p)$ is the kernel of the determinant map to the $p-1$ non-zero elements.
\[ |SL_2| = \frac{(p^2-1)p(p-1)}{p-1} = p(p^2-1) = p(p-1)(p+1). \]

\textbf{Part (2): Orders of elements.}
The characteristic polynomial is $x^2 - \tau x + 1 = 0$ (since $\det=1$).
- \textbf{Case 1: Irreducible over $\mathbb{F}_p$.} The roots lie in $\mathbb{F}_{p^2}$. The roots satisfy $\lambda^{p+1} = 1$ in the multiplicative group. Thus the order divides $p+1$.
- \textbf{Case 2: Reducible with distinct roots.} The roots are in $\mathbb{F}_p$. Since their product is 1, the order divides $p-1$.
- \textbf{Case 3: Repeated root.} The root must be $1$ or $-1$ (since $\lambda^2=1$).
  If $\lambda=1$, the matrix is conjugate to $\begin{pmatrix} 1 & 1 \\ 0 & 1 \end{pmatrix}$, order $p$.
  If $\lambda=-1$, the matrix is $\begin{pmatrix} -1 & 1 \\ 0 & -1 \end{pmatrix}$, order $2p$.
In all cases, the order divides $(p-1)(p+1) = p^2-1$ or $2p$.
\end{solution}

\begin{exercise}
Irreducible complex representations of $S_4$.
\end{exercise}

\begin{solution}
\textbf{Prerequisites / Core Concepts:}
\begin{itemize}
    \item \textbf{Conjugacy Classes:} There are 5 classes for $S_4$ corresponding to partitions of 4.
    \item \textbf{Sum of Squares:} $d_1^2 + d_2^2 + d_3^2 + d_4^2 + d_5^2 = 24$.
\end{itemize}

\textbf{Step 1: Identify the dimensions.}
The only set of 5 squares that sum to 24 is $1^2, 1^2, 2^2, 3^2, 3^2$.

\textbf{Step 2: Describe the representations.}
1. \textbf{Trivial ($d=1$):} Maps every permutation to 1.
2. \textbf{Sign ($d=1$):} Maps even permutations to 1 and odd permutations to -1.
3. \textbf{2D Irrep:} $S_4$ has a normal subgroup $V_4$. $S_4/V_4 \cong S_3$. The 2D irrep of $S_3$ lifts to $S_4$.
4. \textbf{Standard ($d=3$):} $S_4$ acts on $\mathbb{C}^4$ by permuting coordinates. The subspace where $\sum z_i = 0$ is a 3D irrep.
5. \textbf{Vieri ($d=3$):} This is the tensor product of the Standard representation and the Sign representation.
\end{solution}

\begin{exercise}
Let $F$ be the splitting field of $x^4 - 2$.
\end{exercise}

\begin{solution}
\textbf{Prerequisites / Core Concepts:}
\begin{itemize}
    \item \textbf{Splitting Field:} Smallest field containing all roots.
    \item \textbf{Galois Group:} The group of symmetries of the roots.
\end{itemize}

\textbf{Step 1: Find the roots and the field.}
The roots are $\sqrt[4]{2}, -\sqrt[4]{2}, i\sqrt[4]{2}, -i\sqrt[4]{2}$.
The field is $F = \mathbb{Q}(\sqrt[4]{2}, i)$.
The degree is $[F:\mathbb{Q}] = [\mathbb{Q}(\sqrt[4]{2}, i) : \mathbb{Q}(\sqrt[4]{2})] \times [\mathbb{Q}(\sqrt[4]{2}) : \mathbb{Q}] = 2 \times 4 = 8$.

\textbf{Step 2: Determine the Galois Group.}
The automorphisms are determined by where they send $\sqrt[4]{2}$ and $i$.
Let $\sigma(\sqrt[4]{2}) = i\sqrt[4]{2}$ and $\sigma(i) = i$ (order 4 rotation).
Let $\tau(\sqrt[4]{2}) = \sqrt[4]{2}$ and $\tau(i) = -i$ (reflection).
These generate the Dihedral group $D_8$.

\textbf{Step 3: Subfields.}
By the Fundamental Theorem of Galois Theory, subfields correspond to the 10 subgroups of $D_8$:
- The group $D_8$ itself corresponds to $\mathbb{Q}$.
- The trivial group corresponds to $F$.
- Subgroups of order 2 correspond to fields of degree 4, such as $\mathbb{Q}(\sqrt[4]{2})$ and $\mathbb{Q}(i\sqrt[4]{2})$.
- Subgroups of order 4 correspond to quadratic fields like $\mathbb{Q}(i)$ and $\mathbb{Q}(\sqrt{2})$.
\end{solution}
